\documentclass[11pt,a4paper]{article}
\usepackage[utf8]{inputenc}
\usepackage[T1]{fontenc}
\usepackage[french]{babel}
\usepackage[a4paper,top=2cm,bottom=2.1cm,left=2.05cm,right=2.05cm]{geometry}
\usepackage{amsmath,amssymb}
\usepackage{array}
\usepackage{booktabs}
\usepackage{enumitem}
\usepackage{eurosym}
\usepackage{fancyhdr}
\usepackage{hyperref}
\usepackage{inconsolata}
\usepackage{listings}
\usepackage{microtype}
\usepackage{tabularx}
\usepackage[most]{tcolorbox}
\usepackage{titlesec}
\usepackage{xcolor}

\definecolor{ink}{HTML}{202124}
\definecolor{muted}{HTML}{5F6368}
\definecolor{line}{HTML}{DADCE0}
\definecolor{paper}{HTML}{F8F9FA}
\definecolor{accent}{HTML}{8A1538}
\definecolor{accentsoft}{HTML}{F8E8EE}
\definecolor{green}{HTML}{0B6B57}
\definecolor{greensoft}{HTML}{E6F4EF}
\definecolor{amber}{HTML}{9A5B00}
\definecolor{ambersoft}{HTML}{FFF4D6}
\definecolor{blue}{HTML}{245B8E}
\definecolor{bluesoft}{HTML}{EAF2FA}
\definecolor{violet}{HTML}{5C4D7D}
\definecolor{violetsoft}{HTML}{F1EEF8}

\hypersetup{
  colorlinks=true,
  linkcolor=accent,
  urlcolor=green
}

\setlength{\parindent}{0pt}
\setlength{\parskip}{5pt}
\setlist[itemize]{leftmargin=1.35em,itemsep=2pt,topsep=2pt}
\setlist[enumerate]{leftmargin=1.55em,itemsep=3pt,topsep=3pt}
\setlist[enumerate,1]{label=\textbf{\alph*.}}

\pagestyle{fancy}
\fancyhf{}
\fancyhead[L]{Python pour économistes}
\fancyhead[R]{Université de Lille -- L2 Économie}
\fancyfoot[C]{\thepage}
\renewcommand{\headrulewidth}{0.35pt}
\setlength{\headheight}{14.5pt}

\titleformat{\section}
  {\Large\bfseries\color{accent}}
  {}{0pt}{}
  [\vspace{-0.25em}\color{line}\titlerule]
\titleformat{\subsection}
  {\large\bfseries\color{ink}}
  {}{0pt}{}

\lstdefinelanguage{terminal}{
  morekeywords={ls,cd,pwd,mkdir,python,python3,code,dir,type,echo},
  sensitive=true
}

\lstset{
  basicstyle=\ttfamily\small,
  backgroundcolor=\color{paper},
  frame=single,
  framerule=0.35pt,
  rulecolor=\color{line},
  language=Python,
  keywordstyle=\color{accent}\bfseries,
  commentstyle=\color{muted},
  stringstyle=\color{green},
  showstringspaces=false,
  breaklines=true,
  columns=fullflexible,
  keepspaces=true,
  upquote=true,
  literate=
   {à}{{\`a}}1 {â}{{\^a}}1 {ä}{{\"a}}1
   {ç}{{\c c}}1
   {é}{{\'e}}1 {è}{{\`e}}1 {ê}{{\^e}}1 {ë}{{\"e}}1
   {î}{{\^\i}}1 {ï}{{\"\i}}1
   {ô}{{\^o}}1 {ö}{{\"o}}1
   {ù}{{\`u}}1 {û}{{\^u}}1 {ü}{{\"u}}1
   {À}{{\`A}}1 {Â}{{\^A}}1 {Ä}{{\"A}}1
   {Ç}{{\c C}}1
   {É}{{\'E}}1 {È}{{\`E}}1 {Ê}{{\^E}}1 {Ë}{{\"E}}1
   {Î}{{\^I}}1 {Ï}{{\"I}}1
   {Ô}{{\^O}}1 {Ö}{{\"O}}1
   {Ù}{{\`U}}1 {Û}{{\^U}}1 {Ü}{{\"U}}1
   {œ}{{\oe}}1 {Œ}{{\OE}}1
   {–}{{--}}1 {—}{{---}}1
}

\newcommand{\code}[1]{\texttt{#1}}
\newcommand{\pp}{\,points de pourcentage}

\newcounter{exercise}
\newcounter{extension}
\newcounter{logic}

\newcommand{\workbooktitle}[4]{%
  \setcounter{exercise}{0}%
  \setcounter{extension}{0}%
  \setcounter{logic}{0}%
  \fancyhead[L]{#1 -- #2}%
  {\Large\bfseries\color{ink}#1 -- #2\par}
  \vspace{0.2em}
  {\color{muted}Python pour économistes -- Université de Lille, L2 Économie\par}
  \vspace{0.65em}
  {\color{accent}\rule{\linewidth}{0.8pt}}
  \vspace{0.5em}
  \begin{routebox}{Question fil rouge}
  #3
  \end{routebox}
  \begin{deliverablebox}{Livrables}
  #4
  \end{deliverablebox}
  \begin{methodbox}{Méthode de travail}
  \begin{tabularx}{\linewidth}{@{}>{\bfseries\color{blue}}p{0.18\linewidth}X@{}}
  Lire & identifier les données, la question et le livrable attendu.\\
  Décomposer & écrire les étapes en pseudo-code avant de coder.\\
  Exécuter & lancer le script complet depuis le bon dossier.\\
  Vérifier & contrôler les sorties, les fichiers créés et l’interprétation.
  \end{tabularx}
  \end{methodbox}
}

\newtcolorbox{routebox}[1]{
  enhanced,breakable,
  title={#1},
  colback=paper,
  colframe=line,
  coltitle=ink,
  fonttitle=\bfseries,
  boxrule=0.5pt,
  arc=2pt,
  left=8pt,right=8pt,top=7pt,bottom=7pt,
  borderline west={2.4pt}{0pt}{accent}
}

\newtcolorbox{deliverablebox}[1]{
  enhanced,breakable,
  title={#1},
  colback=greensoft,
  colframe=green,
  coltitle=ink,
  fonttitle=\bfseries,
  boxrule=0.6pt,
  arc=2pt,
  left=8pt,right=8pt,top=7pt,bottom=7pt,
  borderline west={2.4pt}{0pt}{accent}
}


\newtcolorbox{methodbox}[1]{
  enhanced,breakable,
  title={#1},
  colback=bluesoft,
  colframe=blue,
  coltitle=ink,
  fonttitle=\bfseries,
  boxrule=0.55pt,
  arc=2pt,
  left=8pt,right=8pt,top=6pt,bottom=6pt,
  before skip=7pt,after skip=9pt,
  borderline west={2.4pt}{0pt}{blue}
}

\newtcolorbox{pathwaybox}[1]{
  enhanced,breakable,
  title={#1},
  colback=violetsoft,
  colframe=violet,
  coltitle=ink,
  fonttitle=\bfseries,
  boxrule=0.55pt,
  arc=2pt,
  left=8pt,right=8pt,top=6pt,bottom=6pt,
  before skip=8pt,after skip=8pt,
  borderline west={2.4pt}{0pt}{violet}
}

\newenvironment{pseudocodeexercise}[1]{%
  \refstepcounter{logic}%
  \begin{tcolorbox}[
    enhanced,breakable,
    title={Logique \thelogic{} -- #1 \hfill\scriptsize PSEUDO-CODE},
    colback=bluesoft,
    colframe=blue,
    coltitle=ink,
    fonttitle=\bfseries,
    boxrule=0.6pt,
    arc=2pt,
    left=8pt,right=8pt,top=6pt,bottom=6pt,
    before skip=8pt,after skip=8pt,
    borderline west={2.6pt}{0pt}{blue}
  ]%
}{\end{tcolorbox}}

\newenvironment{mandatoryexercise}[1]{%
  \refstepcounter{exercise}%
  \begin{tcolorbox}[
    enhanced,breakable,
    title={Exercice \theexercise{} -- #1 \hfill\scriptsize OBLIGATOIRE},
    colback=white,
    colframe=accent,
    coltitle=ink,
    fonttitle=\bfseries,
    boxrule=0.7pt,
    arc=2pt,
    left=8pt,right=8pt,top=7pt,bottom=7pt,
    before skip=9pt,after skip=9pt,
    borderline west={2.4pt}{0pt}{accent}
  ]%
}{\end{tcolorbox}}

\newenvironment{extensionexercise}[1]{%
  \refstepcounter{extension}%
  \begin{tcolorbox}[
    enhanced,breakable,
    title={Extension \theextension{} -- #1},
    colback=ambersoft,
    colframe=amber,
    coltitle=ink,
    fonttitle=\bfseries,
    boxrule=0.55pt,
    arc=2pt,
    left=8pt,right=8pt,top=7pt,bottom=7pt,
    before skip=8pt,after skip=8pt,
    borderline west={2.4pt}{0pt}{amber}
  ]%
}{\end{tcolorbox}}

\newtcolorbox{checkpointbox}[1]{
  enhanced,breakable,
  title={#1},
  colback=accentsoft,
  colframe=accent,
  coltitle=ink,
  fonttitle=\bfseries,
  boxrule=0.55pt,
  arc=2pt,
  left=8pt,right=8pt,top=7pt,bottom=7pt,
  borderline west={2.4pt}{0pt}{accent}
}


\begin{document}

\workbooktitle
  {TP3}
  {Fonctions et comparaison de groupes}
  {Comment écrire des fonctions Python pour comparer des groupes de pays à partir d'indicateurs OWID ?}
  {Un fichier \code{tp3\_nom\_prenom.py}, un tableau texte de statistiques descriptives simples et une interprétation de cinq lignes.}

\begin{checkpointbox}{Parcours obligatoire}
Les exercices 1 à 7 sont obligatoires. Les notions statistiques obligatoires sont : moyenne, médiane, minimum, maximum et proportion. Les quartiles et l'écart-type restent facultatifs.
\end{checkpointbox}

\begin{checkpointbox}{Source des données}
Les valeurs de PIB par habitant sont des extraits pédagogiques arrondis d'Our World in Data.
\begin{lstlisting}[language=terminal]
https://ourworldindata.org/grapher/gdp-per-capita-worldbank.csv
\end{lstlisting}
\end{checkpointbox}

\begin{pseudocodeexercise}{Automatiser une statistique}
\begin{lstlisting}
# Objectif : comparer deux listes de valeurs
# Entrees : liste_groupe_a, liste_groupe_b
# Definir une fonction moyenne(liste)
# Definir une fonction mediane(liste)
# Appliquer les fonctions aux deux groupes
# Comparer les resultats
# Sortie : tableau + phrase d'interpretation
\end{lstlisting}
\end{pseudocodeexercise}

\section*{A. Données et rappels}

\begin{mandatoryexercise}{Créer deux groupes de pays}
Créer deux listes de PIB par habitant, exprimées en milliers de dollars internationaux et arrondies pour l'exercice.
\begin{lstlisting}
groupe_europe = [46.0, 54.0, 34.0, 38.0, 22.0, 61.0]
groupe_monde = [46.0, 76.0, 19.0, 10.0, 7.0, 3.0]
\end{lstlisting}
Ajouter deux listes de noms pour documenter les valeurs :
\begin{lstlisting}
pays_europe = ["France", "Germany", "Spain", "Italy", "Poland", "Sweden"]
pays_monde = ["France", "United States", "China", "Brazil", "India", "Ethiopia"]
\end{lstlisting}
\end{mandatoryexercise}

\begin{mandatoryexercise}{Rappel intuitif}
Ajouter dans le script trois commentaires :
\begin{enumerate}
  \item la moyenne résume un niveau général ;
  \item la médiane est la valeur centrale après tri ;
  \item min et max donnent l'étendue des valeurs observées.
\end{enumerate}
\end{mandatoryexercise}

\section*{B. Fonctions}

\begin{mandatoryexercise}{Fonction moyenne}
Écrire une fonction \code{moyenne(valeurs)} qui renvoie la moyenne d'une liste. La tester sur les deux groupes.
\begin{lstlisting}
def moyenne(valeurs):
    return sum(valeurs) / len(valeurs)
\end{lstlisting}
Afficher les résultats avec une unité : \code{milliers de dollars par habitant}.
\end{mandatoryexercise}

\begin{mandatoryexercise}{Fonction médiane guidée}
Pour une liste de longueur paire, on peut prendre la moyenne des deux valeurs centrales. Compléter la fonction :
\begin{lstlisting}
def mediane_paire(valeurs):
    valeurs_triees = sorted(valeurs)
    n = len(valeurs_triees)
    milieu_droit = n // 2
    milieu_gauche = milieu_droit - 1
    return (valeurs_triees[milieu_gauche] + valeurs_triees[milieu_droit]) / 2
\end{lstlisting}
Tester la fonction sur les deux groupes.
\end{mandatoryexercise}

\begin{mandatoryexercise}{Min, max et proportion}
Pour chaque groupe, calculer :
\begin{enumerate}
  \item la valeur minimale ;
  \item la valeur maximale ;
  \item la proportion de pays au-dessus de 30 milliers de dollars par habitant.
\end{enumerate}
Conseil : utiliser un compteur et une boucle.
\end{mandatoryexercise}

\section*{C. Comparer et interpréter}

\begin{mandatoryexercise}{Tableau texte}
Afficher un petit tableau texte avec les colonnes : \code{groupe}, \code{moyenne}, \code{mediane}, \code{min}, \code{max}, \code{part\_au\_dessus\_30}. L'objectif est la lisibilité, pas la perfection typographique.
\end{mandatoryexercise}

\begin{mandatoryexercise}{Question économique}
Ajouter cinq commentaires à la fin du script :
\begin{enumerate}
  \item quel groupe a la moyenne la plus élevée ?
  \item quel groupe a la médiane la plus élevée ?
  \item pourquoi moyenne et médiane peuvent-elles raconter deux choses différentes ?
  \item pourquoi le choix des pays influence-t-il fortement le résultat ?
  \item quelle limite voyez-vous à travailler sur seulement six pays par groupe ?
\end{enumerate}
\end{mandatoryexercise}

\section*{D. Entraînement facultatif}

\begin{extensionexercise}{Statistiques supplémentaires sans pression}
Choisir au moins deux questions si le parcours obligatoire est terminé.
\begin{enumerate}
  \item Ajouter un pays dans chaque groupe et vérifier que les fonctions marchent encore.
  \item Calculer l'étendue : \code{max - min}.
  \item Écrire une fonction \code{proportion\_au\_dessus(valeurs, seuil)}.
  \item Tester le seuil de 50 au lieu de 30.
  \item Écrire une fonction \code{resume(valeurs)} qui affiche moyenne, médiane, min et max.
  \item Comparer la moyenne avant et après ajout d'une valeur très élevée.
  \item Calculer un premier quartile avec une méthode simple de votre choix.
  \item Calculer l'écart-type à partir d'une formule fournie par l'enseignant.
  \item Écrire une phrase expliquant pourquoi l'écart-type mesure la dispersion.
  \item Transformer les valeurs en dollars au lieu de milliers de dollars.
\end{enumerate}
\end{extensionexercise}

\end{document}
